
\chapter*{Expos\'e III\\Infinitesimal Extensions}
\addcontentsline{toc}{chapter}{Expos\'e III. Infinitesimal Extensions}
\markboth{EXPOS\'E III. INFINITESIMAL EXTENSIONS}%
  {EXPOS\'E III. INFINITESIMAL EXTENSIONS}

\begin{center}
  \textit{by M. Demazure}
\end{center}

In this expos\'e we work in the following general situation. We have a
scheme \(S\) and a coherent nilpotent ideal \(\mathscr I\) on \(S\).
We denote by \(S_n\) the closed subscheme of \(S\) defined by the ideal
\(\mathscr I^{\,n+1}\), for \(n\geq 0\). In particular, \(S_0\) is
defined by \(\mathscr I\). Since \(\mathscr I\) is nilpotent, \(S_n=S\)
for all sufficiently large \(n\), and the \(S_i\) have the same
underlying topological space. A typical example is the following:
\(S\) is the spectrum of a local Artinian ring \(A\), and \(\mathscr I\)
is the ideal defined by the radical of \(A\), so that \(S_0\) is the
spectrum of the residue field of \(A\).

In this situation one is given certain data over \(S_0\) and seeks data
over \(S\) that lift them, that is, data whose base change from \(S\) to
\(S_0\) recovers the original data. This is done step by step, through
the schemes \(S_n\). At each step we shall define the obstructions that
arise and classify the resulting solutions when they exist.

The passage from \(S_n\) to \(S_{n+1}\) may be generalized as follows.
We have a scheme \(S\) and two ideals \(\mathscr I\supset\mathscr J\)
such that \(\mathscr I\mathscr J=0\). In the preceding case \(S\),
\(\mathscr I\), and \(\mathscr J\) are respectively
\[
 S_{n+1},\qquad
 \mathscr I/\mathscr I^{\,n+2},\qquad
 \mathscr I^{\,n+1}/\mathscr I^{\,n+2}.
\]
We denote by \(S_0\), respectively \(S_{\mathscr J}\), the closed
subscheme of \(S\) defined by \(\mathscr I\), respectively
\(\mathscr J\), and pose a problem of extending data from
\(S_{\mathscr J}\) to \(S\).

SGA~1, Expos\'e III, treated the problems of extending morphisms of
schemes and of extending schemes. Here we consider the problems of
extending group morphisms, group structures, and subgroups.

In \S0 we have collected those results of SGA~1, Expos\'e III, that
will be useful to us, putting them into the form best suited to our
purpose and sparing the reader constant reference to that
expos\'e. \SGARefLink{sga3:iii:section:1}{Section~1} collects computations in group cohomology
that will be used later and have nothing to do with scheme theory.
\SGARefLink{sga3:iii:section:2}{Sections~2} and~3 treat, respectively, the extension of group morphisms
and the extension of group structures. In \S4 we briefly recall the
proof of a result stated in TDTE~IV concerning the extension of
subschemes, and apply it to the problem of extending subgroups. For the
remainder of the Seminar, only the result of \S2 concerning the
extension of group morphisms is indispensable.

The idea of reducing problems of infinitesimal extension to the usual
cohomological computations in group extensions was suggested by
J.~Giraud during the oral expos\'e, whose computations were appreciably
more complicated and less transparent. Unfortunately, this method seems
to work well only for the first two problems studied, and we have been
unable to avoid rather laborious computations for extensions of
subgroups.

To simplify the language, we shall call a \(Y\)-functor, respectively a
\(Y\)-scheme, and so on, a functor, respectively a scheme, and so on,
equipped with a morphism to the functor \(Y\). This extends the
definitions of Expos\'e~I, which concerned only representable \(Y\).

\section*{\SGAInlineRefTarget{sga3:iii:section:0}0. Reminders from SGA 1, Expos\'e III, and miscellaneous remarks}
\addcontentsline{toc}{section}{0. Reminders from SGA 1, Expos\'e III}

We begin with a general definition.

\medskip
\SGARefTarget{sga3:iii:definition:0:1}\noindent\textbf{Definition 0.1.}
\textit{Let \(\mathscr C\) be a category, \(X\) an object of
\(\widehat{\mathscr C}\), and \(G\) a
\(\widehat{\mathscr C}\)-group acting on \(X\). We say that \(X\) is
\emph{formally principal homogeneous} under \(G\) if the following equivalent
conditions hold:}
\begin{enumerate}
\item[(i)] \textit{for every object \(T\) of \(\mathscr C\), the set
  \(X(T)\) is either empty or principal homogeneous under \(G(T)\);}
\item[(ii)] \textit{the morphism of functors
  \[
    G\times X\longrightarrow X\times X,
    \qquad (g,x)\longmapsto(gx,x),
  \]
  is an isomorphism.}
\end{enumerate}

We shall now put the results of SGA~1, III~\S5 into the form most
useful here. Throughout this section we use the following notation. On a
scheme \(S\) we have two quasi-coherent ideals \(\mathscr I\) and
\(\mathscr J\) such that
\[
 \mathscr I\supset\mathscr J,
 \qquad
 \mathscr I\mathscr J=0.
\]
In particular, \(\mathscr J^2=0\). We denote by \(S_0\), respectively
\(S_{\mathscr J}\), the closed subscheme of \(S\) defined by
\(\mathscr I\), respectively \(\mathscr J\). For every \(S\)-functor
\(X\), we systematically denote by \(X_0\) and \(X_{\mathscr J}\) the
functors obtained by base change from \(S\) to \(S_0\) and
\(S_{\mathscr J}\). We use the same notation for a morphism.

\medskip
\noindent\SGAInlineRefTarget{sga3:iii:definition:0-1-1}\textbf{Definition 0.1.1.}
\textit{Let \(X\) be an \(S\)-functor. Define a functor \(X^+\) over
\(S\) by}
\[
 \operatorname{Hom}_S(Y,X^+)
 =\operatorname{Hom}_{S_{\mathscr J}}
    (Y_{\mathscr J},X_{\mathscr J})
 =\operatorname{Hom}_S(Y_{\mathscr J},X)
\]
\textit{for a variable \(S\)-scheme \(Y\).} In the notation of
Expos\'e~II, \S1,
\[
 X^+\simeq\prod_{S_{\mathscr J}/S}X_{\mathscr J}.
\]
The identity morphism of \(X_{\mathscr J}\) defines, by construction, an
\(S\)-morphism
\[
 p_X:X\longrightarrow X^+.
\]
Explicitly,\footnote{Editor's note: the following sentence has been
added, and the next two remarks have been expanded.} for every
\(S\)-scheme \(Y\), the map
\[
 p_X(Y):
 \operatorname{Hom}_S(Y,X)\longrightarrow
 \operatorname{Hom}_S(Y,X^+)
 =\operatorname{Hom}_S(Y_{\mathscr J},X)
\]
is induced by the morphism \(Y_{\mathscr J}\to Y\).

\medskip
\noindent\SGAInlineRefTarget{sga3:iii:remark:0-1-2}\textbf{Remark 0.1.2.}
If \(X\) is an \(S\)-group functor, then \(X_{\mathscr J}\) is an
\(S_{\mathscr J}\)-group functor. The defining formula for \(X^+\)
therefore equips \(X^+\) with an \(S\)-group-functor structure, and the
description above shows that
\[
 p_X:X\longrightarrow X^+
\]
is a morphism of \(S\)-group functors.

\medskip
\noindent\SGAInlineRefTarget{sga3:iii:remark:0-1-3}\textbf{Remark 0.1.3.}
For every \(S\)-group functor \(Y\), one has
\[
 \operatorname{Hom}_{S\text{-}\mathrm{gr}}(Y,X^+)
 =
 \operatorname{Hom}_{S_{\mathscr J}\text{-}\mathrm{gr}}
   (Y_{\mathscr J},X_{\mathscr J}).
\]
Indeed, let \(f\in\operatorname{Hom}_S(Y,X^+)\) correspond to
\(f_{\mathscr J}\in
\operatorname{Hom}_{S_{\mathscr J}}
(Y_{\mathscr J},X_{\mathscr J})\). The condition that \(f\) be an
\(S\)-group morphism is that, for every \(T\to S\) and every
\(y,y'\in Y(T)\),
\[
 f(y\cdot y')=f(y)\cdot f(y'),
\]
which is equivalent to
\[
 f_{\mathscr J}(y_{\mathscr J})
 \cdot f_{\mathscr J}(y'_{\mathscr J})
 =f_{\mathscr J}\bigl((y\cdot y')_{\mathscr J}\bigr).
\]
Since
\((y\cdot y')_{\mathscr J}
=y_{\mathscr J}\cdot y'_{\mathscr J}\), this is precisely the
condition that \(f_{\mathscr J}\) be a group morphism. Applying this to
\(Y=X\), we again see that \(p_X\), which corresponds to
\(\operatorname{id}_{X_{\mathscr J}}\), is a morphism of
\(S\)-group functors.

Return now to the general case and suppose that \(X\) is an \(S\)-scheme.
An \(S\)-morphism from a variable \(S\)-scheme \(Y\) to \(X^+\) is, by
definition, an \(S_{\mathscr J}\)-morphism
\(g_{\mathscr J}:Y_{\mathscr J}\to X_{\mathscr J}\). We define an
\(X^+\)-functor in abelian groups \(L_X\) as follows.

\medskip
\noindent\SGAInlineRefTarget{sga3:iii:scholium:0-1-4}\textbf{Scholium 0.1.4.}\footnote{Editor's note: this scholium
has been added.}
If \(\pi:Y\to S\) is a morphism of schemes and \(M\) is an
\(\mathscr O_S\)-module, we denote the inverse image \(\pi^*M\) by
\[
 M\otimes_{\mathscr O_S}\mathscr O_Y.
\]
If \(\mathscr J\) is an ideal of \(\mathscr O_S\), we denote by
\(\mathscr J\mathscr O_Y\) the ideal of \(\mathscr O_Y\) that is the
image of
\[
 \pi^*(\mathscr J)
 =\mathscr J\otimes_{\mathscr O_S}\mathscr O_Y
 \longrightarrow\mathscr O_Y.
\]
For every morphism of \(S\)-schemes \(f:Z\to Y\), there is an
epimorphism of \(\mathscr O_Z\)-modules
\begin{equation}
 f^*(\mathscr J\mathscr O_Y)\longrightarrow
 \mathscr J\mathscr O_Z.
 \tag{0.1.4}\label{sga3:iii:scholium:0-1-4:equation:0-1-4}
\end{equation}
It follows from the commutative diagram below.

\SGAThreeNativeDiagram[0.70\textwidth]
  {native_diagrams/exp3/diagram_001_exp3_p004_scholie014_tensor_diagram.tex}
  {The tensor-product diagram yielding the epimorphism (\SGARefLink{sga3:iii:scholium:0-1-4:equation:0-1-4}{0.1.4}).}
  {fig:exp3-scholie-014-tensor}
  {LNM 151, p.~86}

\medskip
\noindent\SGAInlineRefTarget{sga3:iii:definition:0-1-5}\textbf{Definition 0.1.5.}
Let \(X\) be an \(S\)-scheme. For every \(X^+\)-scheme \(Y\), given by
a morphism
\[
 g_{\mathscr J}:Y_{\mathscr J}\longrightarrow X_{\mathscr J},
\]
set
\[
 \operatorname{Hom}_{X^+}(Y,L_X)
 =
 \operatorname{Hom}_{\mathscr O_{Y_0}}
 \bigl(g_0^*(\Omega^1_{X_0/S_0}),\mathscr J\mathscr O_Y\bigr),
\]
where \(\Omega^1_{X_0/S_0}\) is the module of relative differentials of
\(X_0\) over \(S_0\), and \(\mathscr J\mathscr O_Y\) is regarded as an
\(\mathscr O_{Y_0}\)-module because it is annihilated by
\(\mathscr I\).

Then \(L_X\) is an \(X^+\)-functor in abelian
groups.\footnote{Editor's note: what follows has been added.} Indeed,
for every \(X^+\)-morphism \(f:Z\to Y\), the functor \(f_0^*\) and the
morphism
\[
 f_0^*(\mathscr J\mathscr O_Y)\longrightarrow
 \mathscr J\mathscr O_Z
\]
of (\SGARefLink{sga3:iii:scholium:0-1-4}{0.1.4}) induce a natural homomorphism \(L_X(f)\):
\[
\begin{aligned}
 &\operatorname{Hom}_{\mathscr O_{Y_0}}
   \bigl(g_0^*(\Omega^1_{X_0/S_0}),\mathscr J\mathscr O_Y\bigr)\\
 &\quad\longrightarrow
 \operatorname{Hom}_{\mathscr O_{Z_0}}
   \bigl(f_0^*g_0^*(\Omega^1_{X_0/S_0}),
         f_0^*(\mathscr J\mathscr O_Y)\bigr)\\
 &\quad\longrightarrow
 \operatorname{Hom}_{\mathscr O_{Z_0}}
   \bigl(f_0^*g_0^*(\Omega^1_{X_0/S_0}),
         \mathscr J\mathscr O_Z\bigr).
\end{aligned}
\]

Finally, \(L_X(f)\) has the following simple local description. The
schemes \(Y\) and \(Y_{\mathscr J}\) have the same underlying
topological space, as do \(V\) and \(V_{\mathscr J}\) for every open
subset \(V\) of \(Y\). Let
\[
 U=\operatorname{Spec}(A),\qquad
 V=\operatorname{Spec}(B),\qquad
 W=\operatorname{Spec}(C)
\]
be affine opens such that \(U\) lies over
\(\operatorname{Spec}(\Lambda)\subset S\),
\(g_{\mathscr J}(V_{\mathscr J})\subset U\), and
\(W\subset f^{-1}(V)\). Write \(I,J\) for the ideals of \(\Lambda\)
corresponding to \(\mathscr I,\mathscr J\). Then \(f\), respectively
\(g_{\mathscr J}\), induces a morphism of \(\Lambda\)-algebras
\[
 \theta:B\longrightarrow C,
 \qquad\text{respectively}\qquad
 \varphi:A\longrightarrow B/JB,
\]
and clearly \(\theta(JB)\subset JC\). On the other hand,
\[
 m\in
 \operatorname{Hom}_{\mathscr O_{Y_0}}
 \bigl(g_0^*(\Omega^1_{X/S}),\mathscr J\mathscr O_Y\bigr)
\]
induces an element \(D\) of
\[
\begin{aligned}
 \operatorname{Hom}_{\mathscr O_{V_0}}
 \bigl(g_0^*(\Omega^1_{U/S}),\mathscr J\mathscr O_V\bigr)
 &=
 \operatorname{Hom}_{B/IB}
 \bigl(\Omega^1_{A/\Lambda}\otimes_A B/IB,JB\bigr)\\
 &=\operatorname{Der}_{\Lambda}(A,JB),
\end{aligned}
\]
and the image of \(L_X(f)(m)\) in
\[
\begin{aligned}
 \operatorname{Hom}_{\mathscr O_{W_0}}
 \bigl(f_0^*g_0^*(\Omega^1_{U/S}),\mathscr J\mathscr O_Z\bigr)
 &=
 \operatorname{Hom}_{C/IC}
 \bigl(\Omega^1_{A/\Lambda}\otimes_A C/IC,JC\bigr)\\
 &=\operatorname{Der}_{\Lambda}(A,JC)
\end{aligned}
\]
is simply \(\theta\circ D\).

\medskip
\noindent\SGAInlineRefTarget{sga3:iii:remark:0-1-6}\textbf{Remark 0.1.6.}\footnote{Editor's note: this remark has
been expanded.}
Let \(f:X\to W\) be an \(S\)-morphism. It induces an \(S\)-morphism
\[
 f^+:X^+\longrightarrow W^+
\]
defined as follows. If \(g\in\operatorname{Hom}_S(Y,X^+)\)
corresponds to an \(S\)-morphism
\(g_{\mathscr J}:Y_{\mathscr J}\to X\), then \(f^+(g)\) is the
element \(f\circ g_{\mathscr J}\) of
\(\operatorname{Hom}_S(Y,W^+)\). The following diagram is commutative.

\SGAThreeNativeDiagram[0.44\textwidth]
  {native_diagrams/exp3/diagram_002_exp3_p005_remark016_naturality_square.tex}
  {Naturality of the morphisms \(p_X\) and \(p_W\).}
  {fig:exp3-remark-016-naturality}
  {Polo--Gille re-edition, Expos\'e III}

